
\subsection{问题分析与模型构建：源荷预测不确定性下日前购电调度}
区别于问题一的全信息确定性典型日调度，问题二无法获取调度日未来时刻的真实负荷与光伏出力。仅允许在每日0时刻确定当日全部时段的合同购电计划$g_t$，日内不允许对合同进行调整。
全部调度决策严格服从因果时序约束，任意$t$时刻的购电、储能动作，仅能够由信息集$I_{d,t}$导出，\textbf{调度日尚未发生的真实负荷、光伏出力不得参与决策计算}；当日真实源荷真值仅在本时隙决策完全确定之后，用于功率可行性裁剪、能量平衡校验与费用结算，该原则与全文统一的计划--执行--结算分离闭环框架保持一致。

模型取消日内储能SOC强制闭环约束，储能荷电状态支持跨日连续传递。仿真设置1月作为预热周期完成储能状态初始化，消除初始SOC对正式仿真的影响；正式统计周期选取2--12月，合计334天。

受负荷、光伏预测误差影响，日前制定的购电计划与真实工况会产生偏差。合同电量不足时触发高价应急购电；合同电量富余则产生闲置弃购电量，同一日内两类现象可以同时发生。本文整体求解流程依次为：候选预测模型对比筛选 $\rightarrow$ 历史残差聚类分析 $\rightarrow$ 基于报童理论确定风险分位并构造保守安全轨迹 $\rightarrow$ 日前线性规划求解购电计划 $\rightarrow$ MPC滚动线性规划调度 $\rightarrow$ 实测真值执行修正 $\rightarrow$ 年度费用结算与指标统计。

日前合同一经确定，全天不可修改。每个10分钟时隙开展储能充放电决策时，严格遵循信息边界约束：动作确定前仅使用信息集$I_{d,t}$包含的历史观测、当前SOC与已发布预测；储能控制动作确定完成后，才读取本时段真实负荷、光伏真值，仅用于执行层修正与结算，真实源荷数据不输入任何优化决策环节。

日内储能调度采用模型预测控制(MPC)作为主执行器，基于未来$H=49$个时段经过风险修正的预测序列开展联合优化，仅执行第一个时隙的控制动作；同时搭建信息口径完全一致的预测驱动贪心策略作为对照方案，两组方案的差异仅在于是否利用未来多时段预测信息开展跨时隙储能耦合优化。

\subsection{基于相似日的候选预测模型对比}
每日0时刻，依托信息集$\mathcal{F}_{d,0}\subset I_{d,0}$，仅使用调度日之前的历史观测数据，筛选历史样本生成负荷、光伏出力的点预测，不引入调度日任何未来真实信息。记第$d$日$t$时段负荷、光伏真值为$L_{d,t},V_{d,t}$，点预测输出记为$\hat L_{d,t},\hat V_{d,t}$。
\begin{align}
\hat L_{d,t} &= f_{\text{sim}}(\mathcal{F}_{d,0}), \label{eq:q2-lpred}\\
\hat V_{d,t} &= f_{\text{sim}}(\mathcal{F}_{d,0}). \label{eq:q2-vpred}
\end{align}
式中$f_{\text{sim}}(\mathcal{F}_{d,0})$代表相似日预测函数，输入全部为0时刻前可获取历史数据。

本文构造两组可解释的统计基线预测方案F1、F2开展对比：
\begin{itemize}
    \item F1：负荷采用7日前同一时段历史值，光伏采用前7天同一时段出力平均值；
    \item F2：负荷采用7、14、21日前相同时段中位数，光伏沿用前7天同时段平均。
\end{itemize}
受可用历史回放数据规模限制，本算例不引入LSTM、Transformer等复杂时序深度学习模型，避免小样本条件下过拟合风险。通过MAE、RMSE指标对比，F1方案净负荷预测误差优于F2，因此后续仿真统一选用F1作为基础预测器。

定义预测残差，表征点预测与真实工况之间的偏离：
\begin{equation}
\begin{cases}
e^L_{d,t}=L_{d,t}-\hat L_{d,t}\\
e^V_{d,t}=V_{d,t}-\hat V_{d,t}
\end{cases}
\end{equation}
其中$e^L_{d,t}$、$e^V_{d,t}$分别代表负荷、光伏在第$d$日$t$时段的预测残差。负荷、光伏预测存在明显尾部误差，如图\ref{fig:q2-safe}所示。仅依靠基础点预测直接编制购电计划会低估极端工况，带来高额应急购电成本，因此需要基于历史残差开展不确定性风险建模。

\begin{figure}[htbp]
\centering
\includegraphics[width=0.75\textwidth]{figures/q2-safe.png}
\caption{2--12月负荷预测值与真实负荷对比，离对角线较远样本代表预测尾部误差}
\label{fig:q2-safe}
\end{figure}

\subsection{报童理论驱动的残差聚类与保守安全轨迹构造}
微网日前购电决策等价于单周期报童优化问题：计划购电$p$，正常合同单位成本$c$，电量缺口触发5倍价格的应急购电。目标期望成本可以写为：
\[
J(p)=cp+5c\mathbb E\left[(N-p)^+\right]
\]
$N$为净负荷随机不确定量，$(N-p)^+=\max(N-p,0)$代表电量缺口。由报童模型临界分位条件，理论最优分位水平满足
\[
\alpha^*=1-\frac{c}{5c}=0.8.
\]
该分位由题目给定应急购电的价格惩罚倍数导出，并非人为选取的经验安全系数。

对历史净负荷残差向量开展K‑means聚类，提取典型残差场景；考虑有限历史样本以及储能跨时段耦合作用，仿真主配置选取历史窗口长度$W=35$，实际经验分位取$\alpha=0.805$，与理论值0.8存在小幅偏移。基于该分位叠加安全裕度，由点预测生成保守安全轨迹$\widetilde L_{d,t}$、$\widetilde V_{d,t}$：
\begin{align}
\widetilde L_{d,t} &= \hat L_{d,t}+\Delta L^{\text{q}},\\
\widetilde V_{d,t} &= \hat V_{d,t}+\Delta V^{\text{q}}
\end{align}
式中$\Delta L^{\text{q}},\Delta V^{\text{q}}$为分位对应的安全裕度。

保守安全轨迹仅作为日前优化模型与滚动MPC的输入；时序仿真结算阶段仍然代入真实$L_{d,t},V_{d,t}$完成能量平衡与费用计算，保守轨迹不等同于调度日未来真实工况。

\subsection{日前优化模型与MPC滚动仿真执行规则}
复用第三章通用微网物理模型，结合本问题场景做限定：每日零点求解得到当日完整合同购电序列$\{g_t\}$，日内全部时隙$g_t$保持固定不可修改；允许高价应急购电$u_t>0$；取消每日SOC强制闭环约束。

模型主目标为最小化当日购电费用。有限时域滚动优化容易出现储能末端无代价放空的短视行为，引入储能终端价值罚项用于求解引导：
\begin{equation}
\min \quad \sum_{t\in\text{day}d} c_t g_t + \lambda \cdot S_{\text{end}}
\label{eq:q2obj}
\end{equation}
\textbf{重要说明：}$\lambda \cdot S_{\text{end}}$仅作为优化求解引导罚项，不计入年度实际费用统计结算。

约束复用第三章能量平衡、SOC递推、设备运行边界，叠加本问题专属约束：
\begin{align}
g_t = g_{d,0},\quad & \forall\, t\in \text{当日全部10分钟时段},\\
g_t\ge 0,\quad & u_t\ge 0.
\end{align}

得到日初日前购电计划之后，进入逐时隙MPC滚动仿真流程。每个时隙基于信息集$I_{d,t}$与风险修正后的预测序列开展滚动优化，MPC优化时域$H=49$，求解得到储能充放电控制动作；确定动作之后，代入本时段真实负荷、光伏真值，开展可行性裁剪，更新储能SOC，核算闲置弃购电量$w_t$、应急购电量$u_t$并完成费用结算。时隙优化引入极小罚项，抑制求解器生成储能同时充放电的数值伪解。
> 执行逻辑严格遵循计划--执行--结算分离：MPC优化只访问$I_{d,t}$，未来真实源荷不参与优化；真实值仅用于执行修正与事后结算。

\subsection{年度仿真结果与基线对比分析}
完整年度仿真完成之后，2--12月统计周期的关键汇总指标如表\ref{tab:q2_result_sum}所示。
\begin{table}[htbp]
\centering
\caption{问题二年度运行汇总结果（2--12月）}
\label{tab:q2_result_sum}
\begin{tabular}{lr}
\toprule
指标 & 数值 \\
\midrule
总计划购电量 & 22382650.65 $\ \mathrm{kWh}$ \\
总应急购电量 & 245483.29 $\ \mathrm{kWh}$ \\
应急购电总费用 & 1120245.61 元 \\
总闲置合同电量 & 1596110.94 $\ \mathrm{kWh}$ \\
总弃光量 & 1741589.34 $\ \mathrm{kWh}$ \\
储能总充电量 & 7112133.95 $\ \mathrm{kWh}$ \\
储能总放电量 & 5760889.78 $\ \mathrm{kWh}$ \\
年度合计总费用 & 14786101.35 元 \\
\bottomrule
\end{tabular}
\end{table}

仿真全过程储能荷电状态始终满足$1200\le S_t \le 10800\ \mathrm{kWh}$，未出现边界越界。
图\ref{fig:q2-diagnostics}为全年净负荷热力图，净负荷表现出显著的季节与日内波动特征，仅仅依靠年度平均工况无法构造可行的日前购电方案。

\begin{figure}[htbp]
\centering
\includegraphics[width=0.85\textwidth]{figures/q2-diagnostics.png}
\caption{全年净负荷的季节与日内分布特征}
\label{fig:q2-diagnostics}
\end{figure}

图\ref{fig:q2-days}展示统计周期内每日运行总费用与应急购电量时序变化，可以观察到应急购电风险集中出现在部分工况恶劣的日期。

\begin{figure}[htbp]
\centering
\includegraphics[width=0.85\textwidth]{figures/q2-days.png}
\caption{问题二334日运行费用与应急购电量时序变化}
\label{fig:q2-days}
\end{figure}

为验证本文报童理论驱动保守安全轨迹策略的运行效果，设置两组参照基准方案开展对照仿真：理想信息参照（调度日0时刻完全获知真实源荷，仅作为事后不可实现的理论标杆）、无储能基线。量化结果如表\ref{tab:q2_baseline_compare}所示。

\begin{table}[htbp]
\centering
\caption{不同参照方案年度运行指标对比（2--12月）}
\label{tab:q2_baseline_compare}
\begin{tabular}{lccc}
\toprule
方案 & 年度总费用/元 & 应急购电费用/元 & 应急购电量/kWh \\
\midrule
本文主配置(F1-W35-$\alpha$=0.805-H49) & 14786101.35 & 1120245.61 & 245483.29 \\
理想信息参照（事后基准，不可实施） & 12291024.77 & 0 & 0 \\
无储能参与基线 & 16407319.63 & 0 & 0 \\
\bottomrule
\end{tabular}
\end{table}

理想信息参照作为事后理论下界，可以度量全部源荷预测不确定性带来的整体经济损失；无储能基线不存在储能削峰填谷能力，模型被迫设置更加保守的日前购电规模，因此不会产生高价应急购电，但整体年度费用更高。该结果说明储能带来的经济收益建立在合理预测与风险裕度设计之上。本文基于报童理论的保守安全轨迹策略，在保障储能参与削峰填谷的前提下，通过分位裕度对冲预测尾部风险，以适度增加闲置合同电量为代价压缩高价应急购电规模。

风险裕度存在工程折中关系：裕度设置过高抬升合同采购成本；裕度设置过低则应急购电风险上升，需要结合历史残差统计选取折中参数。仿真观测到部分日期同时出现闲置弃购$w_t>0$以及应急购电$u_t>0$，该现象来源于日内源荷条件的时序波动，部分时段合同富余产生弃购，部分时段供给缺口触发应急购电，两类工况可以在同一日内共存。

在统一信息边界、预测参数保持不变的条件下，对比MPC(H=49)与预测驱动贪心两种执行方案。MPC利用未来多时段预测开展跨储能耦合联合优化，贪心策略仅基于当前时隙信息决策。仿真得到两组方案年度总成本差异很小，说明在当前参数配置下，MPC跨时段调度带来的潜在收益大部分被源荷预测不确定性带来的动作偏差抵消。
